\documentclass[11pt,a4paper,oneside]{book}
\usepackage{mathscommun}
\setniveau{Terminale}
\renewcommand{\ReferenceProgramme}{Spécialité (2026--2027 et 2027--2028) et mathématiques expertes}
\renewcommand{\EditionDocument}{Édition du 24 octobre 2026}
\hypersetup{pdftitle={Terminale - Corrigés expliqués}}
\begin{document}
\titrecollection{Corrigés expliqués}{Spécialité, mathématiques expertes et fondements pour le supérieur.}{Un parcours démontré : maîtriser les définitions, choisir les méthodes, contrôler les hypothèses et préparer la licence avant les spécialisations de master.}
\frontmatter
\chapter*{Mode d'emploi}
Ce volume suit le parcours de terminale générale avec spécialité mathématiques et option mathématiques expertes. Il ne constitue pas un cours des voies technologiques ni de l'option mathématiques complémentaires. La table de couverture du guide distingue les textes applicables en 2026--2027 et à partir de 2027--2028.

Les chapitres 1 à 24 forment le parcours de spécialité, avec les sections signalées en approfondissement, notamment les sphères et les barycentres du chapitre 8. Les chapitres 25 à 35 couvrent l'option expertes. Le chapitre 36 est une passerelle : il donne les fondements de plusieurs preuves et ne prétend pas remplacer une licence.

Les rappels de première sont repris au point d'usage. Pour une remise à niveau plus longue, revenir aux chapitres correspondants du volume de première : raisonnement (1), suites (4--7), dérivation (8--11), exponentielle (12), trigonométrie (14--15), produit scalaire (16), probabilités (19--23).

Lire une première fois pour comprendre les objets ; refaire ensuite les preuves sans support ; enfin chercher les exercices et les synthèses avant de consulter les corrigés. Chaque chapitre possède six exercices numérotés de la même façon dans les deux recueils. Les indications de difficulté distinguent les tâches de base, de raisonnement et d'approfondissement.

Les fondements admis sont recensés dans le guide. Les quantificateurs, la complétude et les notations abstraites supplémentaires sont explicités pour préparer le supérieur ; leur maîtrise complète dépasse les exigences scolaires.
\tableofcontents
\mainmatter
\chapter{Langage, fonctions et calcul : les bases}
\section*{Exercice 1.1 -- Nier une affirmation}
La négation est : il existe un réel $x$ tel que, pour tout entier $n$, $n\le x$. L'affirmation initiale est la propriété d'Archimède, vraie dans $\R$ (chapitre 36). Changer l'ordre des quantificateurs changerait le sens.

\section*{Exercice 1.2 -- Domaine avant calcul}
Le logarithme impose $x>1$ et la racine au dénominateur impose $x<4$. Domaine : $]1,4[$. Sur cet intervalle, le dénominateur est non nul, donc $f(x)=0\Leftrightarrow\ln(x-1)=0\Leftrightarrow x=2$.

\section*{Exercice 1.3 -- Une équivalence à corriger}
Il faut $x\ge0$. Le carré donne $x^2-x-6=(x-3)(x+2)=0$. La racine $-2$ ne vérifie pas le signe initial ; seul $3$ convient. Le carré conserve une égalité, mais sa réciproque demande deux membres de même signe.

\section*{Exercice 1.4 -- Composer}
$(g\circ f)(x)=x^2+1$ et $(f\circ g)(x)=(x+1)^2=x^2+2x+1$. Elles coïncident seulement en $x=0$, et ne sont donc pas égales comme fonctions. La composition n'est pas commutative.

\section*{Exercice 1.5 -- Bijection et restriction}
Sur $[0,+\infty[$ vers $[0,+\infty[$, chaque $y\ge0$ a un unique antécédent positif ou nul $\sqrt y$. La réciproque est donc la racine carrée. Sur $\R$, les nombres opposés non nuls ont la même image, empêchant l'injectivité.

\section*{Exercice 1.6 -- Égalité d’ensembles}
Pour $x\in E$, appartenir au membre de gauche signifie $f(x)\in A\cup B$, donc $f(x)\in A$ ou $f(x)\in B$. Cela équivaut à l'appartenance au membre de droite. L'équivalence pour tout $x$ prouve les deux inclusions.

\chapter{Récurrence et sommes}
\section*{Exercice 2.1 -- Somme impaire}
Au rang $1$, $1=1^2$. Si la somme vaut $n^2$, ajouter $2n+1$ donne $(n+1)^2$. Initialisation et hérédité établissent la formule. On peut aussi sommer $2k-1$ en utilisant la somme des entiers.

\section*{Exercice 2.2 -- Somme géométrique}
Factoriser $3^2$ donne $9(1+3+3^2+3^3+3^4)=9(3^5-1)/2=1089$. Il y a cinq termes ; utiliser six termes produirait une erreur de borne.

\section*{Exercice 2.3 -- Télescopage}
Chaque terme vaut $1/(k+1)-1/(k+2)$. Après annulations, il reste $1/2-1/(n+2)$. Vérifier au rang $1$ donne $1/6$, ce qui confirme les bornes.

\section*{Exercice 2.4 -- Une récurrence fausse}
L'hérédité est vraie mais l'initialisation est fausse : $u_0=2$. La vraie formule est $u_n=2^{2^n}$, vérifiée au rang $0$ puis en élevant au carré. Une implication peut être vraie sans que son hypothèse soit réalisée.

\section*{Exercice 2.5 -- Bernoulli appliqué}
Bernoulli avec $x=0{,}01$ et $n=100$ donne $1{,}01^{100}\ge1+100\times0{,}01=2$. C'est une borne, pas une valeur approchée garantie à une précision donnée.

\section*{Exercice 2.6 -- Divisibilité par récurrence}
Au rang $0$, $4^0-1=0$ est divisible par $3$. Si $4^n-1=3k$, alors $4^{n+1}-1=4(4^n-1)+3=3(4k+1)$. Le résultat suit. Le langage des congruences donnera ensuite une preuve plus courte.

\chapter{Dénombrer sans compter deux fois}
\section*{Exercice 3.1 -- Codes}
Codes : $10\times9\times8\times7=5040$. Nombres : $9\times9\times8\times7=4536$, car la première position exclut zéro mais les suivantes peuvent l'utiliser. Les choix successifs se multiplient.

\section*{Exercice 3.2 -- Comité}
Les deux places restantes forment une partie de taille $2$ des huit autres personnes : $\binom82=28$. L'ordre dans le comité n'intervient pas.

\section*{Exercice 3.3 -- Au moins un}
Compter tous les groupes puis enlever ceux formés uniquement d'hommes : $\binom93-\binom43=84-4=80$. Les deux catégories sont complémentaires et le comptage ne répète aucun groupe.

\section*{Exercice 3.4 -- Réunion}
Exactement une : $(22-9)+(18-9)=22$. Au moins une : $22+18-9=31$. Aucune : $40-31=9$. Retirer une seule fois l'intersection dans la réunion, deux fois dans « exactement une ».

\section*{Exercice 3.5 -- Étoiles et séparateurs}
Les six étoiles et deux séparateurs donnent $\binom82=28$ triplets. Les zones vides correspondent aux composantes nulles, donc doivent être autorisées ; l'ordre des trois composantes compte.

\section*{Exercice 3.6 -- Sous-ensembles}
Choisir librement une partie des $n-1$ autres éléments puis ajouter l'élément fixé donne une bijection avec leurs $2^{n-1}$ parties. La moitié de toutes les parties contient donc cet élément.

\chapter{Coefficients binomiaux et formule du binôme}
\section*{Exercice 4.1 -- Pascal}
Depuis $(1,4,6,4,1)$, on obtient $(1,5,10,10,5,1)$. Sa somme vaut $32=2^5$. Les extrémités valent $1$ et chaque valeur intérieure somme deux voisines de la ligne précédente.

\section*{Exercice 4.2 -- Coefficient}
Le terme de degré deux est $\binom52(-2x)^2=40x^2$. Le signe est positif, car l'exposant de $-2$ est pair.

\section*{Exercice 4.3 -- Somme alternée}
Fixer un élément et l'ajouter ou le retirer à chaque partie échange taille paire et taille impaire. Cette opération est sa propre inverse ; elle donne une bijection entre les deux catégories, donc leurs nombres sont égaux.

\section*{Exercice 4.4 -- Élément distingué}
Compter les couples $(B,b)$ avec $B$ une partie et $b\in B$. Par taille de $B$, on obtient le membre de gauche. En choisissant d'abord $b$ parmi $n$ éléments puis les autres membres librement, on obtient $n2^{n-1}$.

\section*{Exercice 4.5 -- Somme de carrés}
Partager $2n$ objets en deux groupes de $n$. Pour choisir $n$ objets dont $k$ dans le premier groupe, on dispose de $\binom nk\binom n{n-k}=\binom nk^2$ choix. Les valeurs de $k$ partitionnent tous les choix.

\section*{Exercice 4.6 -- Binôme et matrices}
La distributivité donne $A^2+AB+BA+B^2$. La formule $A^2+2AB+B^2$ est valide si $AB=BA$. Sans commutativité, les deux termes croisés doivent rester séparés.

\chapter{Vecteurs, bases et repères de l’espace}
\section*{Exercice 5.1 -- Cube}
$E=(0,0,1)$, $C=(1,1,0)$, donc $\overrightarrow{EC}=(1,1,-1)$. Avec $B=(1,0,0)$ et $H=(0,1,1)$, le milieu est $(1/2,1/2,1/2)$, centre du cube.

\section*{Exercice 5.2 -- Combinaison}
$au+bv+cw=(a+b,b,c)$, donc $b=1$, $a=2$, $c=2$. La décomposition est $2u+v+2w$ ; on vérifie en additionnant les coordonnées.

\section*{Exercice 5.3 -- Plan ou droite}
Non : si $v=2u$ et $u\ne0$, les points sont $A+(s+2t)u$, donc forment une droite. Si les deux directions sont nulles, il ne reste qu'un point. La non-colinéarité est essentielle.

\section*{Exercice 5.4 -- Coplanarité}
Le troisième vaut deux fois le premier plus trois fois le deuxième. Ils sont donc coplanaires et ne constituent pas une base de l'espace.

\section*{Exercice 5.5 -- Droites gauches}
Elles sont disjointes puisque les coordonnées $z$ valent $0$ et $1$ ; leurs directions $(1,0,0)$ et $(0,1,0)$ ne sont pas colinéaires. Deux droites coplanaires non parallèles se rencontreraient ; elles sont donc gauches.

\section*{Exercice 5.6 -- Unicité}
Cette relation donnerait deux écritures du vecteur nul : coefficients $(0,0,0)$ et $(1,2,-3)$. L'unicité des coordonnées dans une base les rend incompatibles. La relation exprimerait aussi $w=(u+2v)/3$, donc la coplanarité.

\chapter{Produit scalaire, orthogonalité et projections}
\section*{Exercice 6.1 -- Angle}
Le produit scalaire vaut $1$, les deux normes valent $\sqrt2$. Ainsi $\cos\theta=1/2$ et $\theta=\pi/3$. Le résultat est un angle non orienté entre $0$ et $\pi$.

\section*{Exercice 6.2 -- Projection sur une droite}
Le coefficient de projection est $2/2=1$, donc $H=(1,1,0)$. $MH^2=1+1+1=3$, d'où distance $\sqrt3$. Le produit $\overrightarrow{HM}\cdot(1,1,0)=0$ vérifie le calcul.

\section*{Exercice 6.3 -- Normalité}
Les deux produits scalaires valent zéro. Les deux directions ne sont pas colinéaires ; toute combinaison est donc orthogonale à $(1,1,1)$ par bilinéarité. Cela prouve la normalité au plan entier.

\section*{Exercice 6.4 -- Deux plans}
Le produit scalaire des normales vaut $2-2+0=0$. Les normales étant non nulles, les plans sont perpendiculaires, quelles que soient leurs constantes dans les équations cartésiennes.

\section*{Exercice 6.5 -- Minimum}
$MH$ est orthogonal au plan, donc à $HK$. Pythagore donne $MK^2=MH^2+HK^2>MH^2$ puisque $HK>0$. Les distances sont positives ou nulles ; prendre les racines donne la stricte inégalité.

\section*{Exercice 6.6 -- Inégalité triangulaire}
Développer le carré : $\|u+v\|^2=\|u\|^2+2u\cdot v+\|v\|^2\le(\|u\|+\|v\|)^2$. Les deux membres comparés sont positifs ou nuls ; la racine carrée conserve l'ordre.

\chapter{Droites, plans et systèmes linéaires}
\section*{Exercice 7.1 -- Équation de plan}
$2(x-1)-y+3(z-2)=0$, soit $2x-y+3z-8=0$. Substituer le point donne zéro ; les coefficients de $x,y,z$ restituent la normale.

\section*{Exercice 7.2 -- Intersection}
La substitution donne $t+1+t+2-t=3+t=6$, donc $t=3$. L'unique intersection est $(3,4,-1)$ ; la somme des coordonnées vaut bien $6$.

\section*{Exercice 7.3 -- Projection sur un plan}
Le défaut vaut $6$, la norme au carré de $(1,1,1)$ vaut $3$. Retirer $2(1,1,1)$ donne $H=(-1,0,1)$. La distance vaut $6/\sqrt3=2\sqrt3$ ; la somme des coordonnées de $H$ est nulle.

\section*{Exercice 7.4 -- Intersection de plans}
La différence impose $2y=-2$, donc $y=-1$ ; la première équation donne $x+z=2$. En posant $z=t$, on obtient la droite $(2-t,-1,t)$, de directeur $(-1,0,1)$.

\section*{Exercice 7.5 -- Paramètres indépendants}
L'égalité des coordonnées impose $s=0$, puis $t=1$. Elles se coupent en $(1,0,0)$. Exiger $s=t$ aurait artificiellement rendu le système incompatible et conduit à une conclusion fausse.

\section*{Exercice 7.6 -- Dégénérescence}
La première égalité est satisfaite par tout point de l'espace ; la seconde par aucun. Aucune ne définit un plan. La condition « normale non nulle » ne doit pas être omise.

\chapter{Sphères, barycentres et mesures}
\section*{Exercice 8.1 -- Sphère et plan}
Le projeté du centre est $(0,0,3)$, à distance $3$. L'intersection est un cercle de rayon $\sqrt{25-9}=4$ dans ce plan. Son équation y est $x^2+y^2=16$.

Pour le second lieu, $MA^2=MB^2$ donne $x^2+y^2+z^2=(x-2)^2+y^2+z^2$, donc $x=1$. C’est le plan médiateur de $[AB]$, normal à $(1,0,0)$ et passant par le milieu $(1,0,0)$.

\section*{Exercice 8.2 -- Tangence}
Le rayon a pour direction $(0,0,1)$ ; le plan normal à ce vecteur passant par le point est $z=1$. L'équation de la sphère y devient $x^2+y^2=0$, donc le seul point commun est le point de tangence.

\section*{Exercice 8.3 -- Barycentre}
La somme des poids est $3\ne0$ et $G=(2A+B)/3=(1,0,0)$. Le point est plus proche du point portant le poids le plus grand ; cette observation contrôle le sens du résultat.

\section*{Exercice 8.4 -- Leibniz}
Appliquer l'identité de Leibniz avec poids $1,1$ : le terme constant vaut $IA^2+IB^2=2(AB/2)^2=AB^2/2$. La somme est minimale pour $M=I$.

\section*{Exercice 8.5 -- Volume}
La base dans $z=0$ est un triangle rectangle d'aire $2\times3/2=3$. La hauteur vaut $4$ car l'arête verticale est normale à ce plan. Le volume est $3\times4/3=4$ unités de volume.

\section*{Exercice 8.6 -- Sphère circonscrite}
Pour $C=(x,y,z)$, l'égalité des distances à $O$ et $(2,0,0)$ donne $x=1$ par soustraction des carrés. Les deux autres égalités donnent $y=z=1$. Le centre est $(1,1,1)$ et le rayon $\sqrt3$.

\chapter{Comprendre la limite d’une suite}
\section*{Exercice 9.1 -- Définition et rang}
$3/(n+1)<1/100$ équivaut à $n+1>300$, donc $n\ge300$. Le rang $299$ donne l'égalité $0{,}01$ et ne convient pas pour l'inégalité stricte.

\section*{Exercice 9.2 -- Bornée mais divergente}
La sous-suite paire tend vers $1$ et l'impaire vers $-1$. Si la suite entière avait une limite, toutes ses sous-suites auraient cette même limite, par la définition « à partir d'un rang ». Les deux limites différentes excluent la convergence.

\section*{Exercice 9.3 -- Gendarmes}
$-1/n\le\sin n/n\le1/n$. Les deux bornes tendent vers zéro, donc la suite tend vers zéro. Il n'est pas nécessaire que $\sin n$ converge.

\section*{Exercice 9.4 -- Non majorée}
Poser $u_{2n}=2n+1$ et $u_{2n+1}=1$. La suite n'est pas majorée, mais revient à $1$ à des indices arbitrairement grands ; le seuil $2$ n'est donc jamais dépassé durablement.

\section*{Exercice 9.5 -- Ordre et limite}
Prendre $\varepsilon=\ell/2$. Pour les grands indices, $|u_n-\ell|<\ell/2$, donc $u_n>\ell/2>0$. Cela n'impose aucun signe aux premiers termes.

\section*{Exercice 9.6 -- Négation de convergence}
Il existe $\varepsilon>0$ tel que, pour tout $N$, il existe $n\ge N$ avec $|u_n|\ge\varepsilon$. Le même $\varepsilon$ doit convenir à tous les rangs $N$. Dire seulement que des termes sont non nuls ne suffit pas.

\chapter{Opérations sur les limites et suites géométriques}
\section*{Exercice 10.1 -- Terme dominant}
En divisant numérateur et dénominateur par $n^2$, on obtient $(5+2/n)/(1-3/n^2)\to5$. Le dénominateur est non nul pour les grands indices et sa limite vaut $1$.

\section*{Exercice 10.2 -- Conjugué}
Multiplier par le conjugué donne $4n/(\sqrt{n^2+4n}+n)=4/(\sqrt{1+4/n}+1)\to2$. L'expression initiale était une forme indéterminée.

\section*{Exercice 10.3 -- Géométriques}
Les deux premières ont module $0{,}7^n\to0$, donc convergent vers zéro. Pour la troisième, les termes pairs tendent vers $+\infty$ et les impairs vers $-\infty$ ; elle n'a ni limite réelle ni limite infinie unique.

\section*{Exercice 10.4 -- Suite affine}
Le point fixe est $6$. Poser $v_n=u_n-6$ donne $v_{n+1}=v_n/2$, $v_0=-5$, donc $u_n=6-5/2^n$. La suite croît vers $6$.

\section*{Exercice 10.5 -- Indétermination}
$(2/n)/(1/n)=2$, tandis que $(1/n)/(1/n^2)=n\to+\infty$. Les morceaux ont les mêmes limites nulles, ce qui démontre qu'une forme $0/0$ ne détermine aucune réponse.

\section*{Exercice 10.6 -- Somme infinie}
C'est la limite de $S_n=\sum_{k=1}^n3^{-k}=(1/3)(1-3^{-n})/(1-1/3)$. Comme $3^{-n}\to0$, la limite vaut $1/2$. La définition passe d'abord par une somme finie.

\chapter{Suites monotones, points fixes et suites adjacentes}
\section*{Exercice 11.1 -- Suite racine}
Le point fixe positif est $L=(1+\sqrt{13})/2$. L'intervalle $[1,L]$ est stable car la racine est croissante et $f(L)=L$. Sur cet intervalle, $3+x-x^2\ge0$, donc $f(x)\ge x$. La suite croît et est majorée par $L$ ; sa limite vérifie l'équation et vaut $L$.

\section*{Exercice 11.2 -- Erreur contractante}
Il suffit de $2^n\ge800$. Comme $2^9=512$ et $2^{10}=1024$, $n\ge10$ convient. La borne ne prétend pas que l'erreur réelle atteint ce maximum.

\section*{Exercice 11.3 -- Point fixe trompeur}
L'équation $L=-L$ donne $L=0$, mais $u_n=(-1)^n$ ne converge pas. Un passage à la limite supposant une convergence non prouvée ne peut pas établir celle-ci.

\section*{Exercice 11.4 -- Adjacentes}
$a_n$ croît, $b_n$ décroît, $a_n\le b_n$ et $b_n-a_n=2/(n+1)\to0$. Elles convergent donc vers une même limite, ici $1$. À chaque rang, cette limite est dans $[a_n,b_n]$.

\section*{Exercice 11.5 -- Héron pour 3}
$u_1=(2+3/2)/2=7/4$, puis $u_2=(7/4+12/7)/2=97/56$. La preuve générale donne $u_n\ge\sqrt3$, décroissance et limite $\sqrt3$. L'erreur suivante vaut l'erreur actuelle au carré divisée par $2u_n$.

\section*{Exercice 11.6 -- Limite ou majorant}
Non. On conclut seulement à une limite réelle inférieure ou égale à $100$. La suite constante égale à $1$ fournit un contre-exemple. La limite est la borne supérieure de l'ensemble des termes, pas un majorant arbitraire.

\chapter{Limites de fonctions et croissances comparées}
\section*{Exercice 12.1 -- Limites latérales}
À droite, le dénominateur tend vers zéro par valeurs positives : limite $+\infty$. À gauche, il est négatif : limite $-\infty$. La droite $x=2$ est asymptote verticale ; il n'y a pas de limite bilatérale unique.

\section*{Exercice 12.2 -- Signe à l’infini}
L'expression vaut $(|x|/x)\sqrt{1+4/x^2}$. La limite est $1$ à $+\infty$ et $-1$ à $-\infty$. Remplacer $\sqrt{x^2}$ par $x$ aurait perdu le second signe.

\section*{Exercice 12.3 -- Exponentielle dominante}
Poser $t=2x$ : $x^5e^{-2x}=t^5/(32e^t)\to0$ par croissance comparée. L'exposant positif $2x$ au dénominateur domine la puissance.

\section*{Exercice 12.4 -- Asymptote oblique}
La division donne $x+1+3/(x-1)$. Donc $y=x+1$ est asymptote aux deux infinis et $x=1$ verticale. La courbe est au-dessus de l'oblique pour $x>1$, au-dessous pour $x<1$.

\section*{Exercice 12.5 -- Limite et valeur}
Pour $x\ne1$, $f(x)=x+1$, donc la limite est $2$. Elle diffère de $f(1)=0$ : la fonction n'est pas continue. Changer la valeur en $1$ à $2$ donnerait le prolongement continu.

\section*{Exercice 12.6 -- Une oscillation écrasée}
Pour $x\ne0$, $|x\sin(1/x)|\le|x|\to0$. L'encadrement donne la limite, même si le facteur $\sin(1/x)$ n'en possède pas.

\chapter{Continuité, valeurs intermédiaires et équations}
\section*{Exercice 13.1 -- Existence et unicité}
La fonction est continue sur $\R$, sa dérivée $3x^2+2$ est strictement positive. Elle est donc strictement croissante. Ses valeurs en $0$ et $1$ sont $-2$ et $1$ : le TVI assure une racine entre les deux, la croissance son unicité globale.

\section*{Exercice 13.2 -- Dichotomie}
L'erreur est au plus $1/2^{n+1}$. Avec $n=9$, elle vaut au plus $1/1024<10^{-3}$ ; avec $n=8$, la borne $1/512$ ne suffit pas. Ne pas confondre longueur de l'intervalle et erreur du milieu.

\section*{Exercice 13.3 -- Prolongement}
Le quotient est le taux d'accroissement de l'exponentielle en $0$, donc sa limite vaut $\exp'(0)=1$. Définir la valeur en zéro égale à $1$ réalise le prolongement.

\section*{Exercice 13.4 -- TVI sans monotonie}
Elle possède deux solutions $\pm1/2$. La continuité ne garantit pas l'unicité. En outre les valeurs aux deux extrémités sont égales à $1$ ; pour utiliser directement le TVI, séparer les intervalles $[-1,0]$ et $[0,1]$.

\section*{Exercice 13.5 -- Image d’une suite}
La fonction $x\mapsto\sqrt{1+x^2}$ est continue sur $\R$ : polynôme positif puis racine continue. L'image de la suite converge donc vers $\sqrt5$. La justification est la continuité, pas une substitution purement formelle.

\section*{Exercice 13.6 -- Additivité continue}
L'additivité impose $f(q)=3q$ pour tout rationnel $q$. Choisir des rationnels $q_n\to\sqrt2$ ; la continuité donne $f(\sqrt2)=\lim f(q_n)=3\sqrt2$.

\chapter{Dériver, composer et contrôler les variations}
\section*{Exercice 14.1 -- Chaîne}
Les dérivées sont $6xe^{3x^2-1}$ et $2x/(1+x^2)$. Le second domaine est tout $\R$ car $1+x^2>0$. Dans chaque cas, la dérivée intérieure multiplie celle de la fonction extérieure.

\section*{Exercice 14.2 -- Quotient}
La dérivée vaut $(1-x^2)/(1+x^2)^2$. Elle est positive sur $]-1,1[$ et négative à l'extérieur. La fonction décroît, puis croît, puis décroît ; elle a un minimum local $-1/2$ en $-1$ et un maximum local $1/2$ en $1$.

\section*{Exercice 14.3 -- Tangente}
$\ln1=0$ et $\ln'(1)=1$, donc la tangente est $y=x-1$. On a $\ln(1+h)=h+h\varepsilon(h)$, avec $\varepsilon(h)\to0$ : l'erreur est petite devant $|h|$, sans être identiquement nulle.

\section*{Exercice 14.4 -- Dérivée nulle}
$f(x)=x^3$ convient : $f'(0)=0$ mais $f$ est strictement croissante et change de signe autour de zéro. Une annulation de dérivée n'est qu'un candidat pour un extremum intérieur.

\section*{Exercice 14.5 -- Borne d’erreur}
Si $x\ne y$, les accroissements finis donnent $\sin x-\sin y=\cos(c)(x-y)$ pour un $c$ entre eux. Comme $|\cos c|\le1$, la borne suit. Si $x=y$, les deux membres sont nuls.

\section*{Exercice 14.6 -- Optimisation}
Avec un côté $x>0$, l'autre vaut $36/x$. Le périmètre est $P(x)=2x+72/x$, de dérivée $2-72/x^2$. Il décroît jusqu'à $6$ puis croît. Le minimum est $P(6)=24$, atteint pour le carré.

\chapter{Convexité, tangentes et inégalités}
\section*{Exercice 15.1 -- Convexe mais décroissante}
La dérivée est $-e^{-x}<0$, donc la fonction décroît. La dérivée seconde est $e^{-x}>0$, donc elle est convexe sur $\R$. Ces deux propriétés portent sur des objets différents : valeur et pente.

\section*{Exercice 15.2 -- Inflexion}
$f''(x)=12x^2-4$ s'annule en $x=\pm1/\sqrt3$, change de signe à chacun et vaut négatif entre eux. Les deux points d'inflexion ont ordonnées $1/9-2/3=-5/9$.

\section*{Exercice 15.3 -- Tangente logarithmique}
Le logarithme est strictement concave ; sa courbe est sous la tangente en $1$, $y=x-1$. Pour l'égalité, $g(x)=x-1-\ln x$ a dérivée $(x-1)/x$ : son unique minimum est $g(1)=0$.

\section*{Exercice 15.4 -- Milieu}
Appliquer la définition de convexité de l'exponentielle à $t=1/2$. La stricte convexité donne égalité seulement si $a=b$. On retrouve aussi l'inégalité arithmético-géométrique pour les nombres $e^a,e^b$.

\section*{Exercice 15.5 -- Un faux critère}
Pour $f(x)=x^4$, $f''(x)=12x^2\ge0$ et $f''(0)=0$. La convexité ne change pas de part et d'autre de zéro ; il n'y a pas d'inflexion.

\section*{Exercice 15.6 -- Lorenz}
La somme vaut $10$. Les points sont $(0,0),(1/4,1/10),(1/2,1/5),(3/4,2/5),(1,1)$. Les pentes successives valent $0{,}4,0{,}4,0{,}8,2{,}4$, croissantes, donc la ligne polygonale est convexe.

\chapter{Logarithme, puissances et modèles multiplicatifs}
\section*{Exercice 16.1 -- Domaine et produit}
Le domaine est $x>2$. L'équation devient $x(x-2)=3$, soit $(x-3)(x+1)=0$. Seul $x=3$ est admissible ; $-1$ vient d'une équation algébrique sur un domaine trop large.

\section*{Exercice 16.2 -- Seuil}
En prenant les logarithmes et divisant par $\ln0{,}8<0$, $n>\ln0{,}1/\ln0{,}8\simeq10{,}32$. Le premier entier est $11$. On peut contrôler $0{,}8^{10}>0{,}1$ et $0{,}8^{11}<0{,}1$.

\section*{Exercice 16.3 -- Limite en zéro}
Poser $t=-\ln x\to+\infty$ : l'expression vaut $-te^{-2t}\to0$ par croissance comparée. Le signe reste négatif près de zéro, ce qui est compatible avec une limite nulle.

\section*{Exercice 16.4 -- Puissance réelle}
$x^{\sqrt2}=e^{\sqrt2\ln x}$. La composition donne une dérivée $e^{\sqrt2\ln x}\sqrt2/x=\sqrt2 x^{\sqrt2-1}$. La base positive est indispensable pour cette définition.

\section*{Exercice 16.5 -- Logarithme d’une somme}
Le domaine est $x>0$. Le membre de droite vaut $\ln(2x)$, donc par injectivité $x+1=2x$, d'où $x=1$. On ne sépare jamais $\ln(x+1)$ en deux logarithmes.

\section*{Exercice 16.6 -- Une limite exponentielle}
Le logarithme de l'expression vaut $n\ln(1+2/n)=2\,[\ln(1+2/n)/(2/n)]\to2$ par la dérivée de $\ln$ en $1$. La continuité de l'exponentielle donne la limite $e^2$.

\chapter{Sinus, cosinus et oscillations}
\section*{Exercice 17.1 -- Angles et signes}
Par périodicité et symétrie, les valeurs sont $1/2$ et $-1/2$. L'équation a exactement les solutions $-\pi/3$ et $\pi/3$ sur l'intervalle demandé. Sur $\R$, il faudrait ajouter les multiples de $2\pi$.

\section*{Exercice 17.2 -- Inéquation}
Le cosinus est pair et décroît de $1$ à $-1$ sur $[0,\pi]$. Il atteint $-1/2$ en $2\pi/3$. L'ensemble des solutions est donc $[-2\pi/3,2\pi/3]$, bornes incluses.

\section*{Exercice 17.3 -- Composition}
La première dérivée est $6x\cos(3x^2)$. Pour la seconde fonction, $f'(x)=-2\sin(2x)$, négative sur $]0,\pi/2[$ puis positive sur $]\pi/2,\pi[$. Elle descend de $2$ à $0$ puis remonte à $2$.

\section*{Exercice 17.4 -- Rectangle inscrit}
L'aire est $4R^2\cos t\sin t=2R^2\sin(2t)$. Elle est au plus $2R^2$, atteinte pour $2t=\pi/2$, donc $t=\pi/4$. Les deux côtés sont alors égaux : le rectangle optimal est un carré de côté $R\sqrt2$. L'aire a bien l'unité d'une longueur au carré.

\section*{Exercice 17.5 -- Limite trigonométrique}
La première limite est $5\lim\sin(5x)/(5x)=5$. Pour la seconde, $1-\cos x=2\sin^2(x/2)$, donc le quotient vaut $\frac12[\sin(x/2)/(x/2)]^2\to1/2$. Les formules s'appliquent pour $x\ne0$ puis on passe à la limite.

\section*{Exercice 17.6 -- Tangente et domaine}
$\tan x=\sin x/\cos x$ est définie hors des $\pi/2+k\pi$. Sa dérivée vaut $1/\cos^2x>0$ sur chaque intervalle de définition. Le domaine n'est pas un intervalle : par exemple $\pi/4<3\pi/4$ mais les tangentes valent $1>-1$. La croissance est une propriété sur chaque composante.

\chapter{Primitives et équations différentielles}
\section*{Exercice 18.1 -- Primitives et domaines}
Sur les positifs : $x^3-2\ln x+C$. Sur les négatifs : $x^3-2\ln(-x)+D$, car la dérivée de $\ln(-x)$ vaut $1/x$. Les constantes sont indépendantes sur ces deux intervalles disjoints.

\section*{Exercice 18.2 -- Condition initiale}
Le point d'équilibre est $3$. Poser $z=y-3$ donne $z'=2z$, donc $z=Ce^{2x}$. La condition initiale impose $C=2$ : $y=3+2e^{2x}$. La substitution vérifie à la fois l'équation et la valeur initiale.

\section*{Exercice 18.3 -- Second membre}
Pour $y_p=ax+b$, l'équation impose $a=ax+b+x$, donc $a=-1$, $b=-1$. La différence $y-y_p$ vérifie $z'=z$. Toutes les solutions sont $y(x)=Ce^x-x-1$.

\section*{Exercice 18.4 -- Refroidissement fictif}
$T(t)=20+60e^{-0{,}2t}$. L'équation $60e^{-0{,}2t}=10$ donne $t=5\ln6\simeq8{,}96$ minutes. Le coefficient $0{,}2$ a l'unité minute inverse ; le modèle suppose une température ambiante constante et un échange proportionnel à l'écart.

\section*{Exercice 18.5 -- Euler instable}
Euler donne $y_{n+1}=(1-4)y_n=-3y_n$, donc $y_n=(-3)^n$ diverge en oscillant. La solution exacte $e^{-4t}$ tend vers zéro. Pour que l'erreur homogène décroisse, il faut $|1-4h|<1$, soit $0<h<1/2$.

\section*{Exercice 18.6 -- Logistique}
En notant $q=Ae^{-rt}$, $y'=Krq/(1+q)^2$. D'autre part $ry(1-y/K)=r[K/(1+q)][q/(1+q)]$, même expression. Comme $q\to0$, $y\to K$ ; $K$ est une capacité limite dans ce modèle. Ce calcul vérifie une famille de solutions ; il ne prouve pas à lui seul l'unicité de toutes les solutions.

\chapter{Intégrale et théorème fondamental de l’analyse}
\section*{Exercice 19.1 -- Aire signée}
L'intégrale vaut $[x^2/2]_{-1}^1=0$. L'aire vaut $\int_{-1}^0(-x)\,dx+\int_0^1x\,dx=1$. Les deux triangles ont aires $1/2$ ; leur compensation concerne l'intégrale, pas l'aire géométrique.

\section*{Exercice 19.2 -- Borne variable}
Poser $G(u)=\int_1^ue^{-t^2}\,dt$. L'intégrande est continue, donc $G'(u)=e^{-u^2}$. Par composition, $F'(x)=2xe^{-x^4}$. Aucune primitive élémentaire explicite n'est nécessaire.

\section*{Exercice 19.3 -- Valeur moyenne}
La moyenne vaut $(1/3)\int_0^3t^2\,dt=3$. La valeur est atteinte en $t=\sqrt3\in[0,3]$. Ce n'est ni la valeur au milieu $9/4$, ni la moyenne des valeurs extrêmes $9/2$.

\section*{Exercice 19.4 -- Positivité stricte}
Supposons $f(c)>0$. Par continuité, $f$ reste au moins $f(c)/2$ sur un petit intervalle de longueur positive contenu dans $[a,b]$ (unilatéral si $c$ est une extrémité). Son intégrale sur cet intervalle est strictement positive, et le reste contribue positivement ou par zéro : contradiction.

\section*{Exercice 19.5 -- Orientation}
Chasles donne $\int_2^5f=1-5=-4$. Le changement d'orientation donne $\int_5^2f=4$. Une intégrale négative n'implique pas que la fonction soit négative en chaque point.

\section*{Exercice 19.6 -- Comparaison}
$e^{-1}\le e^{-x^2}\le1$ donne $e^{-1}\le I\le1$. La tangente à l'exponentielle donne $e^{-x^2}\ge1-x^2$, donc $I\ge2/3$, meilleur que $e^{-1}$. La comparaison porte sur un intervalle orienté de gauche à droite.

\chapter{Calcul intégral, encadrements et suites d’intégrales}
\section*{Exercice 20.1 -- Intégration par parties}
Avec $u=x$, $v'=e^x$, $u'=1$, $v=e^x$, l'intégration par parties donne $[xe^x]_0^1-\int_0^1e^x\,dx=e-(e-1)=1$. Les fonctions choisies sont continûment dérivables sur le segment.

\section*{Exercice 20.2 -- Aire entre deux courbes}
Sur $[0,1]$, $x-x^2=x(1-x)\ge0$. L'aire vaut donc $\int_0^1(x-x^2)\,dx=1/2-1/3=1/6$ unité d'aire. Inverser les fonctions aurait donné l'opposé de l'aire.

\section*{Exercice 20.3 -- Suite d’intégrales}
$0\le I_n\le\int_0^1x^n\,dx=1/(n+1)\to0$. Une intégration par parties avec $u=x^n$, $v=-e^{-x}$ donne $I_n=[-x^ne^{-x}]_0^1+n\int_0^1x^{n-1}e^{-x}\,dx$, d'où la récurrence. La borne prouve la limite sans itérer une formule numériquement instable.

\section*{Exercice 20.4 -- Somme de Riemann}
C'est une somme de rectangles à droite pour $f(x)=1/(1+x)$ sur $[0,1]$. La continuité garantit la convergence vers $\int_0^1f(x)\,dx=[\ln(1+x)]_0^1=\ln2$. Les rectangles à gauche donnent un encadrement dont l'écart vaut $1/(2n)$.

\section*{Exercice 20.5 -- Trapèzes et convexité}
La courbe est sous la corde : $f(x)\le f(a)+(x-a)(f(b)-f(a))/(b-a)$. Intégrer cette fonction affine donne exactement l'aire algébrique du trapèze, $(b-a)(f(a)+f(b))/2$. L'argument reste valable si certaines valeurs sont négatives.

\section*{Exercice 20.6 -- Série harmonique}
Sur $[k,k+1]$, $1/x\le1/k$. Donc $\ln(n+1)=\int_1^{n+1}dx/x\le H_n$. Le membre de gauche tend vers $+\infty$, donc $H_n$ aussi. Pourtant le terme $1/n$ tend vers zéro : cette condition ne suffit pas à la convergence d'une série.

\chapter{Expériences aléatoires et loi binomiale}
\section*{Exercice 21.1 -- Reconnaître le modèle}
Non : les tirages ne sont pas indépendants et la proportion d'as change selon les résultats précédents. Avec remise et mélange, le nombre d'as suivrait $\mathcal B(5,1/8)$. Le modèle sans remise relève de la loi hypergéométrique, étudiée plus tard.

\section*{Exercice 21.2 -- Probabilité exacte}
$P(X=2)=\binom52(1/3)^2(2/3)^3=80/243$. Par complément, $P(X\ge1)=1-(2/3)^5=211/243$. Le calcul par complément évite quatre termes supplémentaires.

\section*{Exercice 21.3 -- Atteindre une fiabilité}
Il faut $1-0{,}8^n\ge0{,}95$, soit $n\ge\ln(0{,}05)/\ln(0{,}8)\simeq13{,}43$. Le plus petit entier est $14$. L'indépendance justifie la puissance représentant les échecs successifs.

\section*{Exercice 21.4 -- Intervalle et quantiles discrets}
Par symétrie, $P(2\le X\le8)=1-2[\binom{10}0+\binom{10}1]/2^{10}=501/512\simeq0{,}9785$. Pour le quantile, $P(X\le7)=968/1024=0{,}9453125<0{,}95$, tandis que $P(X\le8)=1013/1024>0{,}95$. Ainsi $b=8$. La masse discrète empêche en général d'obtenir exactement $0{,}95$.

\section*{Exercice 21.5 -- Indépendance par paires}
Sur les quatre couples équiprobables, chaque événement a probabilité $1/2$. Chaque intersection de deux vaut $\{(1,1)\}$ ou un ensemble singleton, donc a probabilité $1/4=(1/2)^2$. L'intersection des trois vaut $\{(1,1)\}$, de probabilité $1/4\ne1/8$. L'indépendance mutuelle exige aussi les intersections de trois événements.

\section*{Exercice 21.6 -- Optimiser un paramètre}
Le coefficient binomial est constant. Sur $]0,1[$, le logarithme de $p^k(1-p)^{n-k}$ a dérivée $k/p-(n-k)/(1-p)=(k-np)/(p(1-p))$. Il croît jusqu'à $p=k/n$ puis décroît. Aux extrémités, la probabilité vaut zéro ; le maximum unique est donc atteint en $k/n$.

\chapter{Espérance, variance et sommes de variables}
\section*{Exercice 22.1 -- Calcul direct}
$E(X)=-1/4+2/4=1/4$. $E(X^2)=1/4+4/4=5/4$, donc $V(X)=5/4-(1/4)^2=19/16$. La variance est positive et exprimée dans l'unité de $X$ au carré.

\section*{Exercice 22.2 -- Transformation affine}
$E(Y)=5-2\times3=-1$. $V(Y)=(-2)^2V(X)=16$ et $\sigma(Y)=4$. L'écart type est multiplié par $|-2|$, jamais par un facteur négatif.

\section*{Exercice 22.3 -- Somme dépendante}
$X+Y=1$ est constante, donc sa variance est nulle. Chaque variable a variance $1/4$, donc la somme des variances vaut $1/2$. Leur covariance vaut $-1/4$ ; le terme $2\operatorname{Cov}(X,Y)$ rétablit l'égalité générale.

\section*{Exercice 22.4 -- Binomiale}
$E(X)=30$, $V(X)=100\times0{,}3\times0{,}7=21$. Pour la fréquence, $V(X/100)=21/10000=0{,}0021$. La division par $100$ divise la variance par $100^2$.

\section*{Exercice 22.5 -- Sans indépendance}
Poser $I_i=1$ si l'élève $i$ oublie, $0$ sinon. Le nombre est $S=\sum I_i$ et la linéarité donne $E(S)=\sum E(I_i)=20/5=4$. Aucune indépendance n'est nécessaire pour cette espérance ; elle serait nécessaire à la formule binomiale de variance.

\section*{Exercice 22.6 -- Moindres carrés}
Écrire $X-a=(X-E(X))+(E(X)-a)$, développer le carré et prendre l'espérance. Le terme croisé disparaît car $E(X-E(X))=0$. Le dernier carré est minimal et nul pour l'unique choix $a=E(X)$. La moyenne est ainsi la meilleure constante pour l'erreur quadratique.

\chapter{Concentration et loi des grands nombres}
\section*{Exercice 23.1 -- Bienaymé--Tchebychev}
L'inégalité donne $9/6^2=1/4$. Il s'agit d'une majoration universelle, non de la probabilité exacte. Elle utilise la variance $9$, pas l'écart type $3$.

\section*{Exercice 23.2 -- Taille d’échantillon}
Comme $p(1-p)\le1/4$, la probabilité est au plus $1/(4n\times0{,}05^2)=100/n$. Il suffit donc de $n\ge10000$. Cette garantie peut être prudente ; elle ne suppose pas une approximation normale.

\section*{Exercice 23.3 -- Rôle de l’indépendance}
La moyenne vaut toujours $X$ et reste à distance $1/2$ de son espérance. Pour $\varepsilon=1/4$, la probabilité d'écart vaut $1$, quel que soit $n$. Les variables ne sont pas indépendantes ; la variance de leur moyenne ne décroît pas.

\section*{Exercice 23.4 -- Convergence en probabilité}
Pour tout $\varepsilon>0$ fixé, $P(|X_n-m|\ge\varepsilon)\le C/(n^2\varepsilon^2)\to0$. C'est exactement la définition de convergence en probabilité vers la constante $m$. Cela ne dit pas que chaque réalisation est toujours proche de $m$.

\section*{Exercice 23.5 -- Intervalle aléatoire}
L'événement « $p$ appartient à l'intervalle » est $|F_n-p|\le\varepsilon$. Son complément est inclus dans $\{|F_n-p|\ge\varepsilon\}$, majoré par la concentration. C'est l'intervalle qui varie avant observation ; $p$ est un paramètre fixe. On peut intersecter l'intervalle avec $[0,1]$.

\section*{Exercice 23.6 -- Attente géométrique}
$T>n$ signifie $n$ échecs : probabilité $(1-p)^n$. La conditionnelle vaut $(1-p)^{m+n}/(1-p)^m=(1-p)^n=P(T>n)$. C'est la propriété sans mémoire de la loi géométrique ; le rang commence ici à $1$.

\chapter{Algorithmes : calculer, prouver et contrôler l’erreur}
\section*{Exercice 24.1 -- Invariant de boucle}
Initialiser \texttt{p=1}, puis pour \texttt{k} de $1$ à $n$ remplacer \texttt{p} par \texttt{p*k}. Après le passage $k$, l'invariant est $p=k!$, établi par récurrence. La boucle possède $n$ passages ; pour $n=0$ elle n'en possède aucun et renvoie bien $0!=1$.

\section*{Exercice 24.2 -- Coût de la dichotomie}
Il faut $3/2^{n+1}\le10^{-6}$, donc $2^{n+1}\ge3\times10^6$. Comme $2^{21}=2097152$ et $2^{22}=4194304$, $n=21$ convient. Le coût croît logarithmiquement avec la précision demandée.

\section*{Exercice 24.3 -- Newton cyclique}
$f(0)=2$, $f'(0)=-2$, donc $x_1=1$. Puis $f(1)=1$, $f'(1)=1$, donc $x_2=0$. Le cycle se répète. Fixer un nombre maximal d'itérations, contrôler la dérivée et signaler l'absence de convergence ; comparer à une dichotomie lorsque ses hypothèses sont disponibles.

\section*{Exercice 24.4 -- Erreur des rectangles}
La fonction décroît. Avec $h=1/n$, $R_n=h\sum_{k=1}^n(1+k/n)^{-1}\le\ln2\le L_n=h\sum_{k=0}^{n-1}(1+k/n)^{-1}$. Les termes se compensent dans $L_n-R_n=h(1-1/2)=1/(2n)$.

\section*{Exercice 24.5 -- Monte-Carlo et erreur}
Chaque indicatrice a paramètre $p=\pi/4$ ; la variance de $\widehat\pi_n$ vaut $16p(1-p)/n\le4/n$. Tchebychev donne $P(|\widehat\pi_n-\pi|\ge\varepsilon)\le4/(n\varepsilon^2)$. Le modèle uniforme dans le carré relève du prolongement continu précisé au cours.

\section*{Exercice 24.6 -- Résidu et erreur}
Non : la racine vaut $1000$, donc l'erreur est $1000$. La dérivée a module $10^{-9}$ ; le contrôle $|x-\alpha|\le|f(x)|/10^{-9}$ donne ici exactement $1000$. La petite échelle de la fonction rend le résidu trompeur sans contrôle de pente.

\chapter{Construire et calculer les nombres complexes}
\section*{Exercice 25.1 -- Calcul algébrique}
Le produit vaut $2+4i-3i-6i^2=8+i$. Multiplier le quotient par le conjugué du dénominateur donne $(1+i)(2+i)/5=(1+3i)/5$. Le dénominateur réel vaut $|2-i|^2=5$.

\section*{Exercice 25.2 -- Conjugué}
Écrire $z=x+i y$ donne $3x-i y=3-i$. L'identification des parties réelle et imaginaire impose $x=1$, $y=1$. La solution unique est $z=1+i$.

\section*{Exercice 25.3 -- Inverse et module}
$z\overline z=|z|^2>0$ ; multiplier $z$ par le quotient indiqué donne $1$. Puis $1=|zz^{-1}|=|z||z^{-1}|$ et la division par $|z|>0$ donne la formule du module.

\section*{Exercice 25.4 -- Binôme complexe}
On peut élever deux fois au carré : $(1+i)^2=2i$, puis $(2i)^2=-4$. Le binôme donne $1+4i+6i^2+4i^3+i^4=1+4i-6-4i+1=-4$.

\section*{Exercice 25.5 -- Absence d’ordre compatible}
Dans un corps ordonné, tout carré est positif ou nul : si $x\ge0$ cela suit de la multiplication ; sinon $-x>0$ et $x^2=(-x)^2>0$. Or $i^2=-1$ serait positif, alors que $1>0$ impose $-1<0$. Cette contradiction interdit un tel ordre.

\section*{Exercice 25.6 -- Inégalité triangulaire}
$|z+w|^2=|z|^2+|w|^2+2\Re(z\overline w)\le|z|^2+|w|^2+2|z\overline w|=(|z|+|w|)^2$. Les deux membres sont positifs, donc on peut prendre les racines. Pour $z=1,w=i$, $\sqrt2<2$.

\chapter{Arguments, exponentielle imaginaire et trigonométrie}
\section*{Exercice 26.1 -- Forme trigonométrique}
Le module vaut $2$ et un argument est $2\pi/3$, car le cosinus vaut $-1/2$ et le sinus $\sqrt3/2$. Ainsi $z=2(\cos(2\pi/3)+i\sin(2\pi/3))=2e^{2i\pi/3}$. Les arguments sont $2\pi/3+2k\pi$.

\section*{Exercice 26.2 -- Puissance}
$\sqrt3+i=2e^{i\pi/6}$. La formule de Moivre donne la puissance $2^6e^{i\pi}=-64$. Les modules se multiplient et les arguments s'additionnent.

\section*{Exercice 26.3 -- Addition trigonométrique}
Développer $(\cos a+i\sin a)(\cos b+i\sin b)$ puis identifier les parties réelle et imaginaire donne $\cos(a+b)=\cos a\cos b-\sin a\sin b$ et $\sin(a+b)=\sin a\cos b+\cos a\sin b$. Dans le cours, le produit exponentiel est lui-même établi géométriquement avant cette réécriture.

\section*{Exercice 26.4 -- Linéarisation}
Avec $2\cos x=e^{i x}+e^{-i x}$, le cube donne $8\cos^3x=e^{3i x}+e^{-3i x}+3(e^{i x}+e^{-i x})=2\cos3x+6\cos x$. Donc $\cos^3x=(\cos3x+3\cos x)/4$.

\section*{Exercice 26.5 -- Somme de phases}
C'est une somme géométrique : $(1-e^{i nt})/(1-e^{i t})$. Si $t$ est multiple de $2\pi$, tous les termes valent $1$ et la somme vaut $n$ ; la formule quotient n'est pas définie et ne doit pas être utilisée directement.

\section*{Exercice 26.6 -- Argument d’un produit}
La somme $3\pi/2$ est un argument, mais elle n'appartient pas à l'intervalle choisi. Soustraire $2\pi$ donne l'argument principal $-\pi/2$. Les égalités d'arguments sont des congruences modulo $2\pi$.

\chapter{Polynômes, racines et factorisation dans le plan complexe}
\section*{Exercice 27.1 -- Second degré réel}
Le discriminant est $4-20=-16$. Les racines sont $(2\pm4i)/2=1\pm2i$. On contrôle leur somme $2$ et leur produit $5$.

\section*{Exercice 27.2 -- Factorisation}
$P(z)=(z-1)(z^2+z+1)$. Les racines sont $1$ et $(-1\pm i\sqrt3)/2$. Elles sont distinctes et forment les trois racines cubiques de l'unité.

\section*{Exercice 27.3 -- Racine double}
La division par $z-1$ donne $z^2+z-2=(z-1)(z+2)$. Donc $P(z)=(z-1)^2(z+2)$ : deux racines distinctes, $1$ de multiplicité $2$ et $-2$ simple. Le degré compte les multiplicités dans une factorisation complète.

\section*{Exercice 27.4 -- Racines conjuguées}
La conjugaison préserve sommes et produits et fixe les coefficients réels. Donc $P(\overline z)=\overline{P(z)}$. Si $P(z)=0$, le membre de droite est nul. La condition sur les coefficients est essentielle : $P(z)=z-i$ ne possède pas $-i$ pour racine.

\section*{Exercice 27.5 -- Nombre de racines}
Un polynôme non nul de degré $d$ possède au plus $d$ racines distinctes : chaque racine se factorise et diminue le degré d'une unité. Si le polynôme proposé était non nul, ses cinq racines contrediraient $d\le4$.

\section*{Exercice 27.6 -- Racine complexe d’un complexe}
L'identification donne $x^2-y^2=3$, $2xy=4$. De plus $x^2+y^2=|z|^2=|3+4i|=5$. Ainsi $x^2=4$, $y^2=1$ et $xy=2>0$. Les solutions sont $2+i$ et $-2-i$.

\chapter{Géométrie complexe et racines de l’unité}
\section*{Exercice 28.1 -- Distances}
Le premier est le cercle de centre d'affixe $1$ et rayon $2$. Pour le second, écrire $z=x+i y$ donne $(x-1)^2+y^2=(x+1)^2+y^2$, donc $x=0$ : l'axe imaginaire, médiatrice du segment d'extrémités $-1,1$.

\section*{Exercice 28.2 -- Rotation}
Elle s'écrit $z'-1=i(z-1)$. C'est la rotation de centre d'affixe $1$ et d'angle $\pi/2$ modulo $2\pi$. Elle conserve les distances car $|i|=1$.

\section*{Exercice 28.3 -- Angle droit}
$(a-c)/(b-c)=(-1-i)/(1-i)=-i$, imaginaire pur non nul. Son argument est $-\pi/2$, donc les vecteurs $\overrightarrow{CB}$ et $\overrightarrow{CA}$ sont perpendiculaires. On peut aussi vérifier $AC^2+BC^2=2+2=4=AB^2$.

\section*{Exercice 28.4 -- Racines quatrièmes}
Les solutions sont $e^{2ki\pi/4}$ pour $k=0,1,2,3$, soit $1,i,-1,-i$. Leur somme vaut zéro. Elles sont les sommets d'un carré centré à l'origine.

\section*{Exercice 28.5 -- Racines d’un nombre non nul}
$-8i=8e^{-i\pi/2}$. Chaque solution a module $2$ et un argument $-\pi/6+2k\pi/3$, $k=0,1,2$. Les solutions sont donc $\sqrt3-i$, $2i$, $-\sqrt3-i$. Leur cube vaut bien le nombre demandé.

\section*{Exercice 28.6 -- Somme des racines de l’unité}
Avec $q=e^{2i\pi/n}\ne1$, la somme est $(1-q^n)/(1-q)=0$. Le barycentre à poids égaux des sommets du polygone régulier a donc affixe zéro, son centre. Pour $n=1$, la somme serait $1$, d'où l'hypothèse.

\chapter{Divisibilité, congruences et algorithme d’Euclide}
\section*{Exercice 29.1 -- Division euclidienne}
$-17=5(-4)+3$, avec $0\le3<5$. Le quotient est $-4$ et le reste $3$. L'écriture $-17=5(-3)-2$ n'est pas une division euclidienne, car le reste est négatif.

\section*{Exercice 29.2 -- Dernier chiffre}
Modulo $10$, les puissances de $7$ suivent $7,9,3,1$ puis recommencent puisque $7^4\equiv1$. Or $2026\equiv2\pmod4$, donc $7^{2026}\equiv7^2\equiv9\pmod{10}$.

\section*{Exercice 29.3 -- Euclide}
$252=198+54$, $198=3\times54+36$, $54=36+18$, $36=2\times18$. Le dernier reste non nul est $18$. L'identité des diviseurs communs à chaque division justifie l'algorithme.

\section*{Exercice 29.4 -- Divisibilité par 9}
Si $N=\sum a_k10^k$, alors $10\equiv1\pmod9$, donc $N\equiv\sum a_k\pmod9$. Ainsi $9$ divise $N$ exactement lorsque $9$ divise la somme de ses chiffres. Le même raisonnement vaut pour $3$.

\section*{Exercice 29.5 -- Simplification interdite}
Non : $x=1,y=4$ donnent $2x=2\equiv8=2y\pmod6$ mais $1\not\equiv4\pmod6$. La simplification conduirait seulement à $x\equiv y\pmod3$. Pour simplifier modulo $6$ sans changer le module, le facteur devrait être premier avec $6$.

\section*{Exercice 29.6 -- Carrés modulo 4}
Un entier pair s'écrit $2a$, carré $4a^2$ ; un impair s'écrit $2a+1$, carré $4a(a+1)+1$. La somme de deux restes dans $\{0,1\}$ vaut $0,1$ ou $2$ modulo $4$, jamais $3$.

\chapter{Bézout, Gauss et équations diophantiennes}
\section*{Exercice 30.1 -- Bézout explicite}
Euclide : $35=22+13$, $22=13+9$, $13=9+4$, $9=2\times4+1$. Remonter donne $1=9-2(13-9)=3(22-13)-2\times13=3\times22-5(35-22)=8\times22-5\times35$. Ainsi $u=-5,v=8$.

\section*{Exercice 30.2 -- Inverse modulaire}
$7\times8=56\equiv1\pmod{11}$, donc multiplier par $8$ donne $x\equiv24\equiv2\pmod{11}$. Toutes les solutions sont $x=2+11k$, $k\in\Z$.

\section*{Exercice 30.3 -- Équation diophantienne}
Diviser par le pgcd $3$ donne $2x+3y=5$. Une solution est $(1,1)$. Toute autre vérifie $2(x-1)=-3(y-1)$ ; Gauss impose $3\mid(x-1)$. Donc $x=1+3k$, $y=1-2k$, $k\in\Z$. La substitution prouve réciproquement que tous ces couples conviennent.

\section*{Exercice 30.4 -- Aucune solution}
Le membre de gauche est divisible par $7$, mais $5$ ne l'est pas. Il n'existe donc aucune solution. Le critère général exige que le pgcd des coefficients divise le second membre.

\section*{Exercice 30.5 -- Gauss et puissances}
Bézout donne $ua+vb=1$. Au carré, $1=a(u^2a+2uvb)+v^2b^2$. Cette identité de Bézout pour $a,b^2$ prouve qu'ils sont premiers entre eux. Autre voie : aucun nombre premier ne peut diviser simultanément $a$ et $b^2$.

\section*{Exercice 30.6 -- Deux congruences}
Écrire $x=2+3k$. Modulo $5$, $3k\equiv1$, donc $k\equiv2$ car l'inverse de $3$ est $2$. Ainsi $x=8+15t$, $t\in\Z$. La substitution vérifie les deux conditions. Le module commun est le produit car $3$ et $5$ sont premiers entre eux.

\chapter{Nombres premiers, Fermat et chiffrement}
\section*{Exercice 31.1 -- Tester un nombre premier}
S'il était composé, il aurait un diviseur premier au plus $\sqrt{97}<10$. Il suffit de tester $2,3,5,7$. Il n'est divisible par aucun : impair, somme des chiffres $16$, dernier chiffre $7$, et $97=7\times13+6$. Il est donc premier.

\section*{Exercice 31.2 -- Factorisation et pgcd}
$360=2^3\times3^2\times5$, $252=2^2\times3^2\times7$. Le pgcd prend les exposants minimaux : $2^2\times3^2=36$. Le ppcm prend les maximaux : $2^3\times3^2\times5\times7=2520$. Le produit des deux vaut $360\times252$.

\section*{Exercice 31.3 -- Fermat}
Comme $7$ est premier et ne divise pas $3$, $3^6\equiv1\pmod7$. Or $100=6\times16+4$, donc $3^{100}\equiv3^4=81\equiv4\pmod7$.

\section*{Exercice 31.4 -- Infinité des nombres premiers}
$N=p_1\cdots p_r+1>1$ possède un diviseur premier $q$. Aucun $p_i$ ne divise $N$, puisque le reste est $1$. Donc $q$ n'est pas dans la liste. On ne prétend pas que $N$ lui-même soit toujours premier.

\section*{Exercice 31.5 -- Chiffrement jouet}
$ed=21\equiv1\pmod{20}$. Le chiffré est $c=2^3\bmod33=8$. Pour déchiffrer : $8^2\equiv31\equiv-2$, $8^4\equiv4$, donc $8^7\equiv4(-2)8=-64\equiv2\pmod{33}$. Les nombres sont volontairement petits et n'offrent aucune sécurité réelle.

\section*{Exercice 31.6 -- Réciproque de Fermat}
$2^{10}=1024=3\times341+1$, donc $(2^{10})^{34}\equiv1$. Pourtant $341$ est composé. La congruence de Fermat est nécessaire lorsque le nombre est premier et premier avec la base, mais cette seule vérification n'est pas suffisante.

\chapter{Matrices, applications linéaires et systèmes}
\section*{Exercice 32.1 -- Produit non commutatif}
$AB=\begin{pmatrix}2&1\\1&1\end{pmatrix}$ et $BA=\begin{pmatrix}1&1\\1&2\end{pmatrix}$. Ils diffèrent ; l'ordre du produit encode l'ordre de composition des applications.

\section*{Exercice 32.2 -- Inverse et système}
La matrice $A=\begin{pmatrix}2&1\\1&1\end{pmatrix}$ a déterminant $1$ et inverse $\begin{pmatrix}1&-1\\-1&2\end{pmatrix}$. Multiplier $\binom53$ donne $\binom21$. La substitution vérifie les deux équations.

\section*{Exercice 32.3 -- Produit nul}
$A=\begin{pmatrix}1&0\\0&0\end{pmatrix}$ et $B=\begin{pmatrix}0&0\\0&1\end{pmatrix}$ ont produit nul. Ainsi $AB=A0$ avec $B\ne0$ : la simplification échoue. Si $A$ est inversible, multiplier par $A^{-1}$ à gauche permet en revanche de conclure $B=C$.

\section*{Exercice 32.4 -- Rotation matricielle}
L'image de $(x,y)$ est $(x\cos t-y\sin t,x\sin t+y\cos t)$. La somme des carrés vaut $x^2(\cos^2t+\sin^2t)+y^2(\sin^2t+\cos^2t)=x^2+y^2$, les termes croisés s'annulant. C'est le même calcul que $|e^{i t}z|=|z|$.

\section*{Exercice 32.5 -- Trois inconnues}
Soustraire les deux premières équations donne $2y=4$, donc $y=2$. Alors $x+z=4$ et $2x-z=-1$. Leur somme donne $3x=3$, d'où $x=1,z=3$. Ces opérations sur les lignes sont réversibles ; la solution est unique et se vérifie directement.

\section*{Exercice 32.6 -- Application linéaire}
Chaque coordonnée de $T$ est une combinaison linéaire : $T(au+bv)=aT(u)+bT(v)$. Sa matrice a pour colonnes $T(1,0)$ et $T(0,1)$ : $\begin{pmatrix}1&2\\3&-1\end{pmatrix}$. En revanche $S(0,0)=(1,0)\ne(0,0)$ ; $S$ est affine mais non linéaire.

\chapter{Puissances de matrices et suites couplées}
\section*{Exercice 33.1 -- Nilpotence}
Au rang zéro, la formule donne $I$. Si elle est vraie au rang $n$, $(I+nN)(I+N)=I+(n+1)N+nN^2=I+(n+1)N$. La récurrence conclut. Le binôme fonctionnerait aussi car $I$ commute avec $N$.

\section*{Exercice 33.2 -- Puissance par diagonalisation}
Ces vecteurs propres ont valeurs propres $3$ et $1$. Avec $P=\begin{pmatrix}1&1\\1&-1\end{pmatrix}$, $P^{-1}=P/2$ et $A=P\operatorname{diag}(3,1)P^{-1}$. Ainsi $A^n=\frac12\begin{pmatrix}3^n+1&3^n-1\\3^n-1&3^n+1\end{pmatrix}$.

\section*{Exercice 33.3 -- Suites couplées}
Le vecteur colonne au rang $n$ est $A^n\binom10$, où $A$ est la matrice de l'exercice précédent. Donc $u_n=(3^n+1)/2$, $v_n=(3^n-1)/2$. On peut aussi étudier la somme, multipliée par $3$, et la différence, constante.

\section*{Exercice 33.4 -- Récurrence affine vectorielle}
Le point fixe est $L=2B=\binom24$. La différence vérifie $U_{n+1}-L=\frac12(U_n-L)$, donc $U_n=L-2^{-n}L=(1-2^{-n})\binom24$. Elle converge vers $L$.

\section*{Exercice 33.5 -- Fibonacci}
$\binom{F_{n+2}}{F_{n+1}}=\begin{pmatrix}1&1\\1&0\end{pmatrix}\binom{F_{n+1}}{F_n}$. Les termes successifs sont $0,1,1,2,3,5,8$, donc $F_6=8$. Les valeurs propres de la matrice résolvent $\lambda^2-\lambda-1=0$, ce qui conduit plus tard à la formule fermée.

\section*{Exercice 33.6 -- Binôme matriciel}
La distributivité donne $A^2+AB+BA+B^2$. Le regroupement en $2AB$ exige $AB=BA$. Les matrices de l'exercice 32.1 fournissent un contre-exemple explicite puisque leurs deux produits diffèrent.

\chapter{Graphes, chemins et matrices d’adjacence}
\section*{Exercice 34.1 -- Degrés}
La somme des degrés vaut $14$. Chaque arête contribue $2$, donc le graphe possède $7$ arêtes. Cette égalité est nécessaire ; elle ne suffit pas à elle seule pour réaliser une liste arbitraire de degrés.

\section*{Exercice 34.2 -- Adjacence}
La matrice est $M=\begin{pmatrix}0&1&0\\1&0&1\\0&1&0\end{pmatrix}$. Son carré est $\begin{pmatrix}1&0&1\\0&2&0\\1&0&1\end{pmatrix}$. Les deux nombres demandés valent $1$ : $A,B,A$ et $A,B,C$. Une marche peut répéter un sommet.

\section*{Exercice 34.3 -- Connexité}
Il n'y a aucune arête entre composantes. Chacune peut être complète : $\binom22+\binom32=1+3=4$ arêtes. La réalisation par un segment et un triangle atteint cette borne.

\section*{Exercice 34.4 -- Sommets de degré impair}
La somme de tous les degrés est paire. Les degrés pairs contribuent une somme paire ; la somme des degrés impairs doit donc être paire. Or une somme d'entiers impairs a la parité de leur nombre. Ce nombre est donc pair.

\section*{Exercice 34.5 -- Chemins et puissances}
Classer chaque marche selon son avant-dernier sommet $k$. Les $n$ premières arêtes forment une marche de $i$ à $k$, comptée par $(M^n)_{ik}$. La dernière arête de $k$ à $j$ est comptée par $M_{kj}$. Les choix se multiplient et les classes disjointes s'additionnent.

\section*{Exercice 34.6 -- Orientation}
À chaque sommet, une seule arête sortante impose la marche. Les positions se répètent tous les trois pas. Comme $2026\equiv1\pmod3$, la marche unique aboutit à $B$ ; le nombre demandé vaut $1$.

\chapter{Chaînes de Markov et distributions invariantes}
\section*{Exercice 35.1 -- Convention de transition}
$\mu_1=\mu_0P=(0{,}8,0{,}2)$, puis $\mu_2=(0{,}8^2+0{,}2\times0{,}3,0{,}8\times0{,}2+0{,}2\times0{,}7)=(0{,}7,0{,}3)$. Les lignes de $P$ donnent les probabilités de départ vers arrivée.

\section*{Exercice 35.2 -- Distribution invariante}
Si $p_n$ est la probabilité du premier état, $p_{n+1}=0{,}8p_n+0{,}3(1-p_n)=0{,}5p_n+0{,}3$. Le point fixe vaut $0{,}6$ et $p_n=0{,}6+(p_0-0{,}6)0{,}5^n$. Donc $\mu_n\to(0{,}6,0{,}4)$ pour toute distribution initiale ; c'est l'unique distribution invariante.

\section*{Exercice 35.3 -- Périodicité}
$(1/2,1/2)$ est invariante. Mais les distributions issues de $(1,0)$ alternent $(1,0),(0,1)$ : elles ne convergent pas. L'existence, même unique, d'une distribution invariante ne suffit pas à garantir la convergence.

\section*{Exercice 35.4 -- Trois états}
Chaque colonne somme à $1$, donc le produit de la distribution uniforme par $P$ est uniforme. Depuis $(1,0,0)$, un pas donne $(1/2,1/2,0)$ puis deux donnent $(1/4,1/2,1/4)$. L'invariance est une équation, pas une simple observation de symétrie.

\section*{Exercice 35.5 -- Probabilité d’un trajet}
La règle des probabilités conditionnelles et la propriété de Markov donnent $P_{12}P_{22}P_{21}=0{,}2\times0{,}7\times0{,}3=0{,}042$. Il s'agit d'un trajet particulier ; la probabilité d'être en état 1 après trois pas somme tous les trajets possibles.

\section*{Exercice 35.6 -- État absorbant}
Pour rester au premier état, il faut y rester à chaque pas, donc probabilité $q^n$. Le premier passage au second état suit une loi géométrique de paramètre $1-q$, sur $\{1,2,\ldots\}$, de moyenne $1/(1-q)$. La formule de moyenne est un prolongement signalé au chapitre 23 ; la limite de distribution est $(0,1)$.

\chapter{Fondements pour la licence : complétude, intégration et abstraction}
\section*{Exercice 36.1 -- Supremum et maximum}
Tous les éléments sont strictement inférieurs à $1$, donc $1$ majore $A$. Si $M<1$, choisir $n>1/(1-M)$ donne $1-1/n>M$ ; aucun tel $M$ ne majore $A$. Le supremum est $1$, mais il n'appartient pas à $A$ : ce n'est pas un maximum.

\section*{Exercice 36.2 -- Une hypothèse indispensable}
Non. $f(x)=1/x$ sur $]0,1]$ est continue et non bornée. $g(x)=x$ sur $]0,1[$ est bornée mais n'atteint ni son infimum $0$ ni son supremum $1$. Le théorème du cours exige un segment fermé et borné.

\section*{Exercice 36.3 -- Convergence uniforme}
Sur $[0,1]$, $\sup|f_n|=1/n\to0$ : convergence uniforme. Sur $\R$, pour tout $n$, choisir $x=n$ donne $|f_n(x)|=1$, donc aucune uniformité vers zéro. Pourtant pour chaque réel $x$ fixé, $x/n\to0$ : convergence ponctuelle.

\section*{Exercice 36.4 -- Noyau d’une dérivation}
Les images des trois vecteurs de base sont $0,1,2X$. Les colonnes de la matrice sont donc $(0,0,0)$, $(1,0,0)$, $(0,2,0)$, soit $\begin{pmatrix}0&1&0\\0&0&2\\0&0&0\end{pmatrix}$. Le noyau est l'ensemble des polynômes constants, car $b+2cX=0$ impose $b=c=0$.

\section*{Exercice 36.5 -- Échange limite–intégrale}
La positivité et l'inégalité triangulaire intégrale donnent $|\int_a^b(f_n-f)|\le\int_a^b|f_n-f|\le(b-a)\sup_{[a,b]}|f_n-f|\to0$. L'uniformité fournit une même borne d'erreur sur tout le segment, précisément ce qui manque à une convergence seulement ponctuelle.

\section*{Exercice 36.6 -- Audit d’une preuve}
Il faut d'abord montrer existence des termes, bornes et convergence. L'intervalle $[0,2]$ est stable par $f(x)=\sqrt{2+x}$. Comme $u_1\ge u_0$ et $f$ est croissante, une récurrence donne $u_{n+1}\ge u_n$. La suite croissante majorée converge. La continuité autorise alors le passage à la limite. $L\ge0$ et $L^2-L-2=0$ imposent $L=2$, la racine $-1$ étant exclue.

\chapter{Problèmes de synthèse}
Les synthèses S1 à S4 et S8 mobilisent la spécialité et ses approfondissements signalés ; S5 à S7 mobilisent aussi les mathématiques expertes.
\section{S1 -- Héron : convergence et efficacité}
\textbf{Stratégie :} commencer par un invariant de positivité, puis utiliser l'ordre avant de résoudre l'équation de limite.
Tous les termes sont positifs par récurrence. Pour $u>0$, $\frac12(u+a/u)-\sqrt a=(u-\sqrt a)^2/(2u)\ge0$ ; chaque terme après le premier majore donc la racine.
Dès $n\ge1$, $u_{n+1}-u_n=(a-u_n^2)/(2u_n)\le0$. La suite décroît et est minorée par $\sqrt a>0$, donc converge vers $L>0$. La continuité sur $]0,+\infty[$ permet $L=(L+a/L)/2$, d'où $L^2=a$ et $L=\sqrt a$.

L'identité du premier paragraphe donne aussi l'erreur exacte demandée. Pour $a=2$, les termes sont $u_1=3/2$, $u_2=17/12$, $u_3=577/408$. Comme $u_n\ge\sqrt a$, on a $a/u_n\le\sqrt a\le u_n$. Ainsi
\[
\frac{816}{577}\le\sqrt2\le\frac{577}{408},\qquad
\frac{577}{408}-\frac{816}{577}=\frac1{235416}.
\]
Le milieu approche la racine avec erreur au plus $1/470832$. Arrêter lorsque $u_n-a/u_n\le2\varepsilon$ certifie une erreur de milieu au plus $\varepsilon$, sous calcul exact ou contrôle d'arrondi.

Newton pour $f(x)=x^2-a$ donne $x-f(x)/f'(x)=(x+a/x)/2$. La positivité et la forme particulière ont permis la preuve ; d'autres fonctions peuvent produire des cycles, comme au chapitre 24. L'erreur au carré explique l'efficacité locale de Héron, mais les flottants finissent par limiter la précision.

\section{S2 -- Géométrie d’un tétraèdre}
Les intercepts donnent $x/2+y/3+z/6=1$, soit $3x+2y+z=6$. Les trois points satisfont cette équation et ne sont pas alignés ; elle définit donc leur plan. Un normal est $n=(3,2,1)$.
La droite normale issue de $O$ est $(x,y,z)=(3t,2t,t)$. Dans le plan, $14t=6$, donc $H=(9/7,6/7,3/7)$ et $OH=6/\sqrt{14}$ cm. Le signe positif du paramètre correspond au plan situé dans la direction du normal choisi.

Le triangle rectangle $OAB$ a aire $2\times3/2=3\,\mathrm{cm}^2$. La hauteur issue de $C$ vaut $6$ cm ; le volume vaut $3\times6/3=6\,\mathrm{cm}^3$. En prenant $ABC$ pour base, $6=\operatorname{aire}(ABC)\,OH/3$, donc l'aire vaut $3\sqrt{14}\,\mathrm{cm}^2$.

La sphère a centre $(0,0,3)$ et rayon $3$ ; développer $x^2+y^2+(z-3)^2=9$ donne l'équation. Dans le plan $z=3$, elle devient $x^2+y^2=9$ : un cercle de centre $(0,0,3)$ et rayon $3$. Les équations établissent les appartenances et les distances ; une perspective peut les suggérer mais déforme angles et longueurs.

\section{S3 -- Logarithme, intégrale et sommes}
La fonction $x\mapsto1/(1+x)$ décroît sur $[0,1]$. Les rectangles à droite minorent son intégrale, ceux à gauche la majorent. Son intégrale vaut $[\ln(1+x)]_0^1=\ln2$. En soustrayant les sommes, les termes intérieurs disparaissent et $L_n-R_n=(1-1/2)/n=1/(2n)$.
Le milieu a donc erreur au plus $1/(4n)$. Choisir $n=2500$ suffit pour $10^{-4}$ ; le programme calcule les deux sommes ou l'une puis ajoute leur écart. L'arrondi flottant doit être distingué de cette erreur théorique.

Comme $0\le1/(1+x)\le1$, $0\le J_n\le\int_0^1x^n=1/(n+1)$. L'identité $x^n/(1+x)+x^{n+1}/(1+x)=x^n$ donne la récurrence. En partant de $J_0=\ln2$, elle conduit par récurrence à
\[
\ln2=\sum_{k=1}^{n}\frac{(-1)^{k-1}}k+(-1)^nJ_n.
\]
Le reste a module au plus $1/(n+1)$, et son signe est connu. La limite de la somme alternée vaut donc $\ln2$. Les rectangles fournissent ici une meilleure borne à nombre de termes comparable, mais d'autres algorithmes, dont Briggs, peuvent converger plus vite. Une impression décimale n'est un certificat que si l'erreur totale est contrôlée.

\section{S4 -- Surbooking : modèle et décision}
$X_n$ est somme de $n$ Bernoulli indépendantes : $X_n\sim\mathcal B(n,0{,}9)$, $E(X_n)=0{,}9n$, $V(X_n)=0{,}09n$. Le risque exact est
$\sum_{k=101}^n\binom nk0{,}9^k0{,}1^{n-k}$, nul si $n=100$.
Partir de $n=100$, calculer le risque pour $n+1$, avancer tant qu'il est au plus $0{,}05$, puis retourner le dernier entier admissible. Un calcul rationnel ou des probabilités numériques stables permettent de vérifier la frontière. Le programme livré trouve $n=106$ ; les risques pour $106$ et $107$ valent environ $0{,}0398$ et $0{,}0806$.

Coupler les variables par $X_{n+1}=X_n+B_{n+1}$ avec $B_{n+1}\ge0$ montre l'inclusion $\{X_n>100\}\subset\{X_{n+1}>100\}$. Pour $0{,}9n>100$, Tchebychev donne $P(X_n\le100)\le0{,}09n/(0{,}9n-100)^2\to0$ ; la recherche finit donc par dépasser le seuil.
Pour $n=105$, $E=94{,}5$, $V=9{,}45$. L'événement $X_n>100$ implique $X_n\ge101$, donc un écart au moins $6{,}5$. La borne est $9{,}45/6{,}5^2\simeq0{,}224$, trop grande pour décider la contrainte. Une borne supérieure au seuil ne prouve pas que le risque le dépasse.

Le modèle suppose une probabilité identique et une indépendance entre personnes. Groupes, météo, annulations communes et estimation de $0{,}9$ peuvent invalider ces hypothèses. Le calcul n'inclut ni les conséquences du refus d'entrée ni les contraintes pratiques : il illustre un raisonnement, sans constituer une politique réelle.

\section{S5 -- Complexes, rotations et polygones}
Si $\omega^j=\omega^k$, alors $(j-k)/5$ est entier ; avec $0\le j,k<5$, cela impose $j=k$. Tous les points sont sur le cercle unité et les angles successifs valent $2\pi/5$ : le polygone est régulier. La somme géométrique est $(1-\omega^5)/(1-\omega)=0$, donc le barycentre est l'origine.

$1-e^{it}=e^{it/2}(e^{-it/2}-e^{it/2})=-2ie^{it/2}\sin(t/2)$ ; prendre les modules donne la formule. Le côté vaut $2\sin(\pi/5)$, la diagonale $2\sin(2\pi/5)$.
Diviser $1+\omega+\omega^2+\omega^3+\omega^4=0$ par $\omega^2$ donne $1+s+(s^2-2)=0$. Or $s=2\cos(2\pi/5)>0$, car $0<2\pi/5<\pi/2$. Donc $s=(-1+\sqrt5)/2$ et $\cos(2\pi/5)=(\sqrt5-1)/4$.

La matrice est $R=\begin{pmatrix}\cos(2\pi/5)&-\sin(2\pi/5)\\\sin(2\pi/5)&\cos(2\pi/5)\end{pmatrix}$. Composer cinq rotations additionne les angles ; $R^5=R_{2\pi}=I$. C'est aussi la traduction de $\omega^5=1$ par l'identification du plan à $\mathbb C$.

\section{S6 -- RSA : preuve et limites}
Euclide donne $40=13\times3+1$, donc $1=40-13\times3$ et $d\equiv-13\equiv27\pmod{40}$. Ainsi $ed=81$.
Le chiffré vaut $7^3=343\equiv13\pmod{55}$. Pour déchiffrer, les carrés successifs donnent $13^2\equiv4$, $13^4\equiv16$, $13^8\equiv36$, $13^{16}\equiv31$. Comme $27=16+8+2+1$, le produit $31\times36\times4\times13$ est congru à $7$ modulo $55$.

Si $5\mid m$, $m^{81}\equiv m\equiv0\pmod5$. Sinon Fermat donne $m^4\equiv1\pmod5$, donc $m^{81}=m(m^4)^{20}\equiv m$. De même modulo $11$ : si $11\mid m$, résultat immédiat ; sinon $m^{10}\equiv1$ et $81=1+8\times10$. Ainsi $5$ et $11$ divisent $m^{81}-m$. Étant premiers entre eux, leur produit divise cette différence par Gauss.

L'exponentiation rapide évite les puissances immenses ; la preuve de correction ne prouve pas la sécurité. Les facteurs minuscules sont immédiatement retrouvés, et le RSA brut sans encodage approprié possède d'autres faiblesses. L'exemple sert à relier congruences, Bézout et Fermat ; les protocoles cryptographiques réels demandent une étude spécialisée.

\section{S7 -- Trois états et convergence}
Les coefficients sont positifs et chaque ligne somme à $1$. Les colonnes somment aussi à $1$, donc $\pi P=\pi$. Chaque coefficient de $J^2$ vaut $3$, donc $J^2=3J$ et $Q^2=Q$. La décomposition annoncée donne des coefficients diagonaux $1/3+(1/4)(2/3)=1/2$ et hors diagonale $1/3-(1/4)(1/3)=1/4$.

On a $Q(I-Q)=(I-Q)Q=0$ et $(I-Q)^2=I-Q$. Par récurrence, les produits mixtes s'annulent et $P^n=Q+4^{-n}(I-Q)$, y compris $n=0$ où la somme vaut $I$.
Toute distribution ligne $\mu_0$ satisfait $\mu_0Q=\pi$ ; ainsi
$\mu_n=\pi+4^{-n}(\mu_0-\pi)\to\pi$.
Si $\nu$ est invariante, $\nu=\nu P^n$ pour tout $n$, donc le passage à la limite donne $\nu=\pi$ : unicité démontrée.

Chaque coordonnée de $\mu_0$ appartient à $[0,1]$, donc son écart à $1/3$ est au plus $2/3$. Il suffit de $(2/3)4^{-n}\le10^{-3}$, soit $4^n\ge2000/3$. Comme $4^4=256$ et $4^5=1024$, $n=5$ convient. Ce contrôle quantifie la convergence ; il ne se limite pas à l'observer sur quelques itérations.

\section{S8 -- Une fonction définie par une intégrale}
L'intégrande est continue, donc le théorème fondamental donne $F'(x)=e^{-x^2}>0$ ; la fonction croît strictement. De plus $F''(x)=-2xe^{-x^2}\le0$ pour $x\ge0$, donc elle est concave.
La tangente à l'exponentielle donne $e^{-u}\ge1-u$. Pour $0\le t\le x\le1$, $1-t^2\le e^{-t^2}\le1$. Intégrer donne $x-x^3/3\le F(x)\le x$.

Pour $x\ge1$, $F(x)=F(1)+\int_1^xe^{-t^2}\,dt\le1+\int_1^xe^{-t}\,dt\le1+e^{-1}$. Une fonction croissante majorée possède une limite finie (prendre le supremum de ses valeurs et utiliser la même preuve que pour les suites). Cette limite est positive, puisque $F(1)\ge2/3$.
Pour $y\ge x\ge1$, $0\le F(y)-F(x)\le\int_x^ye^{-t}\,dt\le e^{-x}$. Faire tendre $y$ vers l'infini donne le contrôle du reste.

Une intégrale définit une fonction même lorsqu'aucune expression avec les fonctions élémentaires usuelles n'est disponible. Nous n'avons pas démontré ici une impossibilité d'expression, ni évalué $L$ : sa valeur $\sqrt\pi/2$ nécessite une méthode supplémentaire, étudiée en licence. Après normalisation, $e^{-t^2}$ intervient dans une densité gaussienne ; il faudra définir probabilités continues et intégrales impropres pour rendre ce lien complet.

\end{document}
